\chapter*{前言}
\addcontentsline{toc}{chapter}{前言}

分析中的无限对象，通常不会一次完整地呈现在我们面前。我们能够掌握的，往往只是有限多个邻域、有限个坐标、有限批测试函数、有限个矩形，或者一列越来越细的分割。由这些局部而有限的信息恢复整体对象，需要同时回答三个问题：观测怎样彼此相容，极限为何存在，以及怎样确认得到的极限正是所寻找的对象。

第一章考察收敛的拓扑语言。序列何时足够，何时必须让索引方式适应有限条件？网与滤子由此出现，并把紧性、乘积拓扑和半连续性置于同一框架。随后，讨论转向连续函数空间：局部紧性、等度连续与函数代数说明，逐点信息在什么条件下能够提升为一致控制。可数性、完备性与选择原则则标明这些结论在一般空间和可分度量空间中各自呈现的形态。

第二章从集合的大小转向可数可加的质量。Riemann 与 Jordan 理论的边界促使我们构造外测度与 Lebesgue 测度；生成类与正则性回答测度如何被识别，推前与乘积说明测度怎样随映射和坐标传输，而分布函数把一维测度化为可以直接计算的增量。章末回到欧氏空间的局部线性化，建立微分、行列式和 Lipschitz 几何之间的联系。

第三章把测度扩张成积分。简单函数提供有限模型，单调收敛和支配收敛负责通过极限；随后，乘积积分、换元与核把积分放回多变量计算。Radon--Nikodym 定理与 $L^p$ 空间则从另一侧回答密度如何存在、函数列怎样完备，以及连续线性观测如何由积分配对表示。章末的收敛模式与向量值积分说明，尖峰、质量逃逸和不可分值域分别需要不同的控制。

第四章沿相反方向从整体平均恢复局部结构。Vitali 覆盖和极大函数导出 Lebesgue 微分；BV 与 Stieltjes 理论把跳跃和奇异增长记入导数测度；mollifier、分布与 Sobolev 空间使粗糙函数仍能被微分；卷积与半群最终把局部平滑、热演化和无界生成元放在同一条实分析链上。

第五章把概率论放回这套分析语言。law 是推前测度，期望是积分，条件期望是对较粗信息的 Radon--Nikodym 表示；独立性、正则条件分布、一致可积与 tightness 分别控制耦合、条件版本、尾部质量和弱极限的存在。集中不等式、大数律、特征函数与中心极限定理则展示，有限次指数或振荡测试怎样识别独立和的尺度。

第六章把单个随机变量扩展为时间索引族。相容有限维 law 经 Kolmogorov 扩张成为过程，矩估计再筛选路径版本；滤过、鞅与停时把动态信息写成条件平均；Markov 核和半群组织状态演化，Poisson 与 Brownian 模型分别呈现跳跃和连续噪声。章末从简单可预测过程的等距完备化建立 \Ito{} 积分，并以 \Ito{} 公式连接随机微分方程和热方程。

第七章以复参数、振荡测试和正泛函重新审视前六章的对象。最小限度的复分析工具使 Stieltjes 变换可以识别测度；Fourier 级数与变换把卷积和半群化为乘法；矩问题显示正性、紧性与唯一性缺一不可；Daniell--Stone 和 Riesz--Markov 则把一致的有限测试重新装配成 Radon 测度。章末的三个综合问题把这些路线汇合起来，同时准确划出通往一般对偶理论、弱星紧性与谱理论仍需跨越的边界。

本书假定读者熟悉本科阶段的实分析、线性代数、度量空间与基础点集拓扑，并了解 Banach space 与 Hilbert space 的基本概念。测度论、概率论、复分析以及更进一步的泛函分析结论均不作为先修。正文中的定义和证明按出现次序展开；读者在遇到抽象概念时，不妨先追问它修补了哪一个具体失败，再比较它在例子和反例中的表现。
